\SourceLeaf{025}\par
Huic solutioni quasdam annotationes adiicere conveniet.

I. Quatenus spectando observationum errores tamquam quantitates primi ordinis, quantitates ordinum altiorum negliguntur, in formula nostra pro \(\lambda,\lambda',\lambda''\) etc. etiam valores eos quotientium \(\frac{\mathrm{d}U}{\mathrm{d}V}\) etc. adoptare licebit, qui prodeunt e valoribus observatis quantitatum \(V,V',V''\) etc. Quoties \(U\) est functio linearis, manifesto nulla prorsus erit differentia.

II. Si loco errorum mediorum observationum, harum pondera introducere malumus, sint haec, secundum unitatem arbitrariam, resp. \(p,p',p''\) etc., atque \(P\) pondus determinationis valoris ipsius \(U\) e valoribus observatis quantitatum \(V,V',V''\) etc. prodeuntis. Ita habebimus
\[
P=\frac{1}{\dfrac{\lambda\lambda}{p}+\dfrac{\lambda'\lambda'}{p'}+\dfrac{\lambda''\lambda''}{p''}+\text{etc.}}
\]

III. Si \(T\) est functio alia data quantitatum \(V,V',V''\) etc. atque, pro harum valoribus veris,
\[
\frac{\mathrm{d}T}{\mathrm{d}V}=\kappa,\quad\frac{\mathrm{d}T}{\mathrm{d}V'}=\kappa',\quad\frac{\mathrm{d}T}{\mathrm{d}V''}=\kappa''\text{ etc.}
\]
error in determinatione valoris ipsius \(T\), e valoribus observatis ipsarum \(V,V',V''\) etc. petita, erit \(=\kappa e+\kappa'e'+\kappa''e''+\text{etc.},\ =E'\), atque error medius in ista determinatione metuendus \(=\BareRoot{(\kappa\kappa mm+\kappa'\kappa'm'm'+\kappa''\kappa''m''m''+\text{etc.})}\). Errores \(E,E'\) vero manifesto ab invicem iam non erunt independentes, valorque medius producti \(EE'\). secus ac valor medius producti \(ee'\), non erit \(=0\), sed \(=\kappa\lambda mm+\kappa'\lambda'm'm'+\kappa''\lambda''m''m''+\text{etc.}\)

IV. Problema nostrum etiam ad casum eum extendere licet, ubi valores quantitatum \(V,V',V''\) etc. non immediate per observationes inveniuntur, sed quomodocunque ex observationum combinationibus derivantur, si modo singularum determinationes ab invicem sunt independentes, i. e. observationibus diversis superstructae: quoties autem haec conditio locum non habet, formula pro \(M\) erronea evaderet. E. g. si una alterave observatio, quae ad determinationem valoris ipsius \(V\) inserviit, etiam ad valorem ipsius \(V'\) determinandum adhibita esset, errores \(e\) et \(e'\) haud amplius ab invicem independentes forent, neque adeo producti \(ee'\) valor medius \(=0\). Si vero in tali casu nexus quantitatum \(V,V'\) cum observationibus simplicibus, e quibus deductae sunt, rite perpenditur, valor \SourceLeaf{026}medius producti \(ee'\) adiumento annotationis III. assignari, atque sic formula pro \(M\) completa reddi poterit.

\Article{19}
Sint \(V,V',V''\) etc. functiones incognitarum \(x,y,z\) etc., multitudo illarum \(=\pi\), multitudo incognitarum \(=\rho\), supponamusque, per observationes vel immediate vel mediate valores functionum inventos esse \(V=L,V'=L',V''=L''\) etc., ita tamen ut hae determinationes ab invicem fuerint independentes. Si \(\rho\) maior est quam \(\pi\), incognitarum evolutio manifesto fit problema indeterminatum; si \(\rho\) ipsi \(\pi\) aequalis est. singulae \(x,y,z\) etc. in formam functionum ipsarum \(V,V',V''\) etc. redigi vel redactae concipi possunt, ita ut ex harum valoribus observatis valores istarum inveniri possint, simulque adiumento art. praec. praecisionem relativam singulis his determinationibus tribuendam assignare liceat; denique si \(\rho\) minor est quam \(\pi\), singulae \(x,y,z\) etc. infinitis modis diversis in formam functionum ipsarum \(V,V',V''\) etc. redigi, adeoque illarum valores infinitis modis diversis erui poterunt. Quae determinationes exacte quidem quadrare deberent, si observationes praecisione absoluta gauderent; quod quum secus se habeat, alii modi alios valores suppeditabunt, nec minus determinationes e combinationibus diversis petitae inaequali praecisione instructae erunt.

Ceterum si in casu secundo vel tertio functiones \(V,V',V''\) etc. ita comparatae essent, ut \(\pi-\rho+1\) ex ipsis, vel plures, tamquam functiones reliquarum spectare liceret, problema respectu posteriorum functionum etiamnum plus quam determinatum esset, respectu incognitarum \(x,y,z\) etc. autem indeterminatum; harum scilicet valores ne tunc quidem determinare liceret, quando valores functionum \(V,V',V''\) etc. absolute exacti dati essent: sed hunc casum a disquisitione nostra excludemus.

Quoties \(V,V',V''\) etc. per se non sunt functiones \textit{lineares} indeterminatarum suarum, hoc efficietur, si loco incognitarum primitivarum introducuntur ipsarum differentiae a valoribus approximatis. quos aliunde cognitos esse supponere licet. Errores medios in determinationibus \(V=L,V'=L',V''=L''\) etc. metuendos resp. denotabimus per \(m,m',m''\) etc., determinationumque pondera per \(p,p',p''\) etc., ita ut sit \(pmm=p'm'm'=p''m''m''\) etc. Rationem, quam inter se tenent errores medii, cognitam supponemus, ita ut pondera, quorum unum ad lubitum accipi potest, sint nota. Denique statuemus
